% !TeX root = ../navier-stokes-manuscript.tex
% ============================================================
% 2.1 Vortex Stretching
% ============================================================

The momentum equation has four pieces:
\[
  \underbrace{\partial_t u}_{\text{local acceleration}}
  +
  \underbrace{(u\cdot\nabla)u}_{\text{nonlinear transport}}
  =
  \underbrace{-\nabla p}_{\text{pressure}}
  +
  \underbrace{\nu\Delta u}_{\text{viscous diffusion}}.
\]
The central regularity problem lies in the competition between nonlinear transport and viscosity.
The hope is simple: viscosity smooths the fluid faster than nonlinear transport can sharpen it.
If that remains true at every scale, the velocity stays smooth.

\subsection{Vortex Stretching}\label{sub:ThePerfectlyStretchingVortex}

Picture one narrow material vortex tube, a rotating portion of fluid carried and deformed by the flow.
Choose a segment along its axis and follow it as the fluid moves.
If \(e\) is its axial direction and \(L(t)\) its length, the symmetric part of the velocity gradient is
\[
  S:=\frac12\bigl(\nabla u+(\nabla u)^{\mathsf T}\bigr),
\]
and the segment is being stretched when \(e\) is an aligned positive direction:
\[
  Se=\sigma(t)e,
  \qquad
  \sigma(t)>0,
  \qquad
  \frac{\dot L(t)}{L(t)}
  =
  e\cdot Se
  =
  \sigma(t).
\]

A positive \(\sigma\) lengthens the segment.
Because the fluid is incompressible, its cross-section contracts at the same time.
Let \(\mathcal V\) be the fixed material volume, set \(A(t):=\mathcal V/L(t)\), and write \(r_{\mathrm{eff}}(t):=\sqrt{A(t)/\pi}\).
Then
\[
  \nabla\cdot u=0,
  \qquad
  \frac{d}{dt}(A L)=0,
  \qquad
  \frac{\dot A}{A}=-\sigma(t),
  \qquad
  \frac{\dot r_{\mathrm{eff}}}{r_{\mathrm{eff}}}
  =
  -\frac{\sigma(t)}{2}.
\]
The effective radius falls at half the fractional rate at which the length grows.
The rotating fluid is becoming longer and thinner in one coupled motion.

Its rotation is measured by the vorticity \(\omega:=\nabla\times u\).
Along the moving fluid it obeys
\[
  \frac{D\omega}{Dt}
  =
  S\omega+\nu\Delta\omega,
  \qquad
  \frac{D}{Dt}
  :=
  \partial_t+u\cdot\nabla,
\]
and when the vorticity is aligned with the stretching direction, \(\omega=|\omega|e\),
\[
  S\omega
  =
  \sigma(t)\omega,
  \qquad
  \frac12\frac{D|\omega|^2}{Dt}
  =
  \sigma(t)|\omega|^2
  +
  \nu\,\omega\cdot\Delta\omega.
\]
The strain contribution now amplifies the aligned vorticity while viscous diffusion remains in the same evolution law.
On an interval where the full right-hand side stays positive, the rotation strengthens.

The stronger rotation appears inside the velocity gradient, while the segment energy still depends only on its speed and fixed volume.
If the speed remains bounded by \(U_0\) on the fixed-volume material segment \(\Omega_{\mathrm{seg}}(t)\), then
\[
  E_{\mathrm{seg}}(t)
  :=
  \frac12\int_{\Omega_{\mathrm{seg}}(t)}|u(x,t)|^2\,dx
  \leq
  \frac12U_0^2\mathcal V
  <
  \infty,
  \qquad
  \left|\frac12\bigl(\nabla u-(\nabla u)^{\mathsf T}\bigr)\right|_{\mathrm F}
  =
  \frac{|\omega|}{\sqrt2}
  \leq
  |\nabla u|_{\mathrm F}.
\]
Because the speed is bounded and the material volume is fixed, the kinetic energy remains finite while the rotating part of the velocity gradient grows.

\divider{\NavierStokes{} Scaling}

The tube now has a smaller cross-section, with its rotating gradient concentrated into the smaller effective radius.
The scale test uses another solution for comparison, not a later state of this material tube.
Fix \(t_0\) and set \(\ell:=r_{\mathrm{eff}}(t_0)\).
For \(\lambda>1\), define
\[
  u_\lambda(x,t)
  :=
  \lambda u(\lambda x,\lambda^2t),
  \qquad
  p_\lambda(x,t)
  :=
  \lambda^2p(\lambda x,\lambda^2t).
\]
At time \(\lambda^{-2}t_0\), the corresponding vortex in \(u_\lambda\) has width \(\lambda^{-1}\ell\).
Differentiating the rescaled fields gives
\[
\begin{aligned}
  \partial_tu_\lambda(x,t)
  &=
  \lambda^3(\partial_tu)(\lambda x,\lambda^2t),
  \\
  (u_\lambda\cdot\nabla)u_\lambda(x,t)
  &=
  \lambda^3\bigl((u\cdot\nabla)u\bigr)(\lambda x,\lambda^2t),
  \\
  \nabla p_\lambda(x,t)
  &=
  \lambda^3(\nabla p)(\lambda x,\lambda^2t),
  \\
  \nu\Delta u_\lambda(x,t)
  &=
  \lambda^3\nu(\Delta u)(\lambda x,\lambda^2t).
\end{aligned}
\]
Local acceleration, nonlinear transport, pressure, and viscosity all gain the same \(\lambda^3\) factor.
The finer width has given viscosity no scaling advantage over nonlinear sharpening.

\divider{Critical Height}

The equation keeps its four-term balance under rescaling, but the Sobolev size of the velocity changes with the height at which it is measured.
Let \(\Lambda:=(-\Delta)^{1/2}\).
A change of variables gives
\[
  \norm{u_\lambda(t)}_{\dot H^s}^2
  =
  \lambda^{2s-1}
  \norm{u(\lambda^2t)}_{\dot H^s}^2.
\]
The exponent vanishes at \(s=\tfrac12\).
Set
\[
  H(t)
  :=
  \frac12\norm{u(t)}_{\dot H^{1/2}}^2
  =
  \frac12\norm{\Lambda^{1/2}u(t)}_2^2.
\]
With \(H_\lambda\) denoting the critical height of \(u_\lambda\), the energy, critical height, and first-derivative energy scale as
\[
  \norm{u_\lambda(t)}_2^2
  =
  \lambda^{-1}\norm{u(\lambda^2t)}_2^2,
  \qquad
  H_\lambda(t)=H(\lambda^2t),
  \qquad
  \norm{\nabla u_\lambda(t)}_2^2
  =
  \lambda\norm{\nabla u(\lambda^2t)}_2^2.
\]
Kinetic energy falls as first derivatives rise.
Between them, the critical height remains fixed.

\divider{Nonlinear Production versus Viscous Dissipation}

The invariant height lets the competition be read along the solution itself.
Apply \(\Lambda^{1/2}\) to the momentum equation and pair with \(\Lambda^{1/2}u\).
Write the signed nonlinear production as
\[
  \mathcal P_{1/2}(t)
  :=
  -\operatorname{Re}
  \pair
    {\Lambda^{1/2}u(t)}
    {\Lambda^{1/2}\bigl((u(t)\cdot\nabla)u(t)\bigr)}_{L^2}.
\]
The pressure pairing vanishes because \(\nabla\cdot u=0\), while the Laplacian contributes \(-\nu\norm{\Lambda^{3/2}u}_2^2\).
The critical height obeys
\[
  H'(t)
  +
  \nu\norm{\Lambda^{3/2}u(t)}_2^2
  =
  \mathcal P_{1/2}(t).
\]
It rises exactly when nonlinear production exceeds viscous dissipation.
Under the rescaling, the two contributions become
\[
  \mathcal P_{1/2}[u_\lambda](t)
  =
  \lambda^2\mathcal P_{1/2}[u](\lambda^2t),
  \qquad
  \nu\norm{\Lambda^{3/2}u_\lambda(t)}_2^2
  =
  \lambda^2\nu\norm{\Lambda^{3/2}u(\lambda^2t)}_2^2.
\]
Neither contribution gains a scaling advantage.
Their common \(\lambda^2\) factor is the inverse of the time contraction \(t\mapsto\lambda^{-2}t\).

\divider{The Heat Clock}

The time contraction has a direct viscous meaning.
For viscosity alone, the linear heat equation \(v_t=\nu\Delta v\) gives
\[
  \widehat v(\xi,t+\delta t)
  =
  e^{-\nu|\xi|^2\delta t}\widehat v(\xi,t).
\]
A structure of width \(r\) is concentrated around frequencies \(|\xi|\simeq r^{-1}\).
The multiplier changes by order one when \(\nu|\xi|^2\delta t\simeq1\), which gives the comparison time
\[
  \tau_\nu(r)
  :=
  \frac{r^2}{\nu}.
\]
This heat clock measures how long viscosity alone needs to change a structure at width \(r\) by order one.
At the finer width,
\[
  \tau_\nu(\lambda^{-1}r)
  =
  \lambda^{-2}\tau_\nu(r),
\]
so the viscous comparison time contracts by the same factor carried by the full \NavierStokes{} rescaling.

\divider{The Zeno Cascade}

Repeating the decrease in width repeats the decrease in its heat clock.
For the geometric sequence
\[
  r_n:=2^{-n}r_0,
  \qquad
  \tau_n:=\frac{r_n^2}{\nu},
\]
the successive clocks satisfy
\[
  \tau_n
  =
  \frac{r_0^2}{\nu}4^{-n},
  \qquad
  \sum_{n=0}^{\infty}\tau_n
  =
  \frac{4r_0^2}{3\nu}
  <
  \infty.
\]
Their sum is finite, so the clocks alone do not prevent infinitely many scale transitions from accumulating before a finite time.
They do not show that a \NavierStokes{} solution performs those transitions.
For this schedule to describe a candidate singularity, one solution would have to pass through widths
\[
  r_0>r_1>r_2>\cdots\longrightarrow0
\]
at times
\[
  t_0<t_1<t_2<\cdots\uparrow T_*<\infty,
  \qquad
  t_{n+1}-t_n\asymp\tau_n,
\]
with comparison constants independent of \(n\).
The remaining time would then satisfy
\[
  T_*-t_n
  \asymp
  \sum_{k=n}^{\infty}\tau_k
  =
  \frac{4r_n^2}{3\nu}.
\]
As \(n\) grows, the same velocity field would have to form each new concentration in less time than the last.
The spatial cascade has therefore become a hierarchy of temporal demands: the velocity must change, its rate of change must itself change, and that rate must change again.
The objects that record those demands are \(u,\partial_tu,\partial_t^2u,\ldots\): the temporal Tower.
